% ============================================================
% Confidence-Bounded Matched Separability
% Rewritten CSA version with direct separability lower bounds.
% ============================================================
\documentclass[10pt,conference]{IEEEtran}
\IEEEoverridecommandlockouts

\usepackage{cite}
\usepackage{amsmath,amssymb,amsfonts}
\usepackage{graphicx}
\usepackage{textcomp}
\usepackage{xcolor}
\def\BibTeX{{\rm B\kern-.05em{\sc i\kern-.025em b}\kern-.08em
    T\kern-.1667em\lower.7ex\hbox{E}\kern-.125emX}}

\usepackage{amsmath,amssymb,amsthm}
\usepackage[T1]{fontenc}
\usepackage[utf8]{inputenc}
\usepackage[expansion=false]{microtype}
\usepackage[margin=1in]{geometry}
\usepackage{booktabs}
\usepackage{tabularx}
\usepackage{longtable}
\usepackage{array}
\usepackage{multirow}
\usepackage{makecell}
\usepackage{colortbl}
\usepackage{hyperref}
\usepackage{algorithm}
\usepackage{algpseudocode}
\usepackage{cleveref}
\usepackage{enumitem}
\usepackage{caption}
\usepackage{subcaption}
\usepackage{mathtools}
\hypersetup{
  colorlinks=true, linkcolor=blue, citecolor=blue, urlcolor=blue
}

\newtheorem{proposition}{Proposition}
\newtheorem{lemma}{Lemma}
\newtheorem{definition}{Definition}
\newtheorem{corollary}{Corollary}
\newtheorem{observation}{Observation}
\theoremstyle{remark}
\newtheorem{remark}{Remark}

% Shorthand
\newcommand{\LCB}{\mathrm{LCB}}
\newcommand{\UCB}{\mathrm{UCB}}
\newcommand{\CSA}{\mathrm{CSA}}
\newcommand{\CSACP}{\mathrm{CSA}_{\mathrm{CP}}}
\newcommand{\CSAW}{\mathrm{CSA}_{\mathrm{W}}}
\newcommand{\thetaP}{\hat{\theta}_P}
\newcommand{\thetaN}{\hat{\theta}_N}
\newcommand{\shat}{\hat{s}_G}
\newcommand{\reasonFac}{r_{\mathrm{sem}}}
\newcommand{\cueFac}{c_{\mathrm{cue}}}
\newcommand{\presentationFac}{\pi_{\mathrm{pres}}}
\newcommand{\EvalSet}{\mathcal{D}_{\mathrm{eval}}}
\newcommand{\CatMetric}{M_{\mathrm{cat}}}
\newcommand{\CatScore}{\operatorname{Score}_{\mathrm{cat}}}
\newcommand{\posarm}{P_G}
\newcommand{\negarm}{N_G}
\newcommand{\tauCSA}{\tau_{\mathrm{CSA}}}
\newcommand{\ADRw}{\mathrm{ADR}^{+w}}
\newcommand{\ADRwu}{\mathrm{ADR}^{+wu}}
\newcommand{\Gmean}{\mathrm{G\text{-}mean}}
\newcommand{\BA}{\mathrm{BA}}
\newcommand{\MCC}{\mathrm{MCC}}
\newcommand{\CSAboot}{\mathrm{CSA}_{\mathrm{boot}}}

\title{%
  \textbf{Supplementary Material:}\\
  Confidence-Bounded Matched Separability
}

\author{
  [Author Names Omitted for Blind Review]
}

\date{}

\begin{document}

\maketitle

\begin{abstract}
This supplementary material gives the formal details, estimator derivations, threshold-boundary definitions, and complete separability-only tables for the confidence-bounded separability ability (CSA) framework.  CSA combines ADR-style matched evaluation~\cite{zhang2026satisfiabilityADR} with classical binomial confidence bounds~\cite{clopper1934use,wilson1927probable,brown2001interval}.  The revision intentionally removes the earlier margin-over-baseline framing.  The reported quantity is the lower-confidence bound on matched positive/negative separability itself.
\end{abstract}

\section{Supplementary Roadmap}

The supplementary material is organized as follows.  \Cref{sec:supp-claim} restates the revised measurement claim and clarifies what CSA can and cannot certify.  \Cref{sec:supp-theory} derives the factorized separability estimator and its confidence-bounded variants.  \Cref{sec:supp-protocol} describes the computation protocol for SAT and vulnerability detection.  \Cref{sec:supp-results} reports complete per-size tables and threshold-parametric boundaries, including the exact numeric values for every point plotted in the main-paper CSA trajectory figures.  \Cref{sec:supp-threats} gives additional validity considerations.

\section{Revised Measurement Claim}
\label{sec:supp-claim}

\subsection{Evaluation Instances, Semantic Factors, and Non-Semantic Cues}

We write an evaluation instance as
\[
x=R(\reasonFac,\cueFac,\presentationFac),\qquad y=g(\reasonFac),
\]
where $\reasonFac$ is the semantic factor intended to determine the label, $\cueFac$ contains non-semantic label-correlated cues, and $\presentationFac$ contains label-preserving presentation variation.  The evaluation design should vary $\reasonFac$ through matched positive/negative contrasts while controlling $\presentationFac$ through canonicalization, randomization, or stability checks.  The cue factor is not subtracted from the CSA score; instead, cue audits are reported separately as diagnostic evidence about possible alternative explanations.

This change is central.  The main metric no longer attempts to produce a margin over an explicit non-semantic reference class.  It reports a conservative lower bound on observed two-sided separability.  That makes the claim narrower but cleaner: CSA is a \emph{separability certificate}, not a mechanism-identification certificate.

\subsection{Revised Core Claim I}

Let $\EvalSet=\{(x_i,y_i)\}_{i=1}^{n}$ and let $\hat y_i=f(x_i)$.  Any ordinary category metric of the form
\[
\CatMetric(f,\EvalSet)
=
\CatScore\bigl((y_1,\ldots,y_n),(\hat y_1,\ldots,\hat y_n)\bigr)
\]
depends only on the ground-truth and predicted label vectors.

\begin{observation}[Label-vector non-identifiability]
If two models produce the same predicted label vector on $\EvalSet$, every category metric of the above form assigns them the same score.
\end{observation}

\begin{proposition}[Evidential insufficiency]
Ordinary category metrics can summarize benchmark agreement, but they cannot by themselves certify reasoning ability, two-sided matched recall, or robustness to non-semantic cue explanations.
\end{proposition}

\begin{proof}
The observation proves invariance to all changes in the model that preserve the predicted labels.  Therefore, a semantic solver and a non-semantic predictor with identical outputs receive identical category scores.  Moreover, a one-sided predictor can score moderately on a balanced benchmark even though one arm of the intended contrast is not recognized.  Additional design information or additional measurements are required for a stronger reasoning claim.
\end{proof}

\subsection{What the Revised Claim Allows}

The revised claim supports the following statement:

\begin{quote}
Given a matched group construction, category predictions, and a confidence protocol, $\CSACP$ lower-bounds the model's observable two-sided separability on that construction.
\end{quote}

It does not support the stronger statement that the model internally used a particular reasoning algorithm.  If witnesses, proofs, executable tests, verified patches, or causal traces are available, they should be evaluated directly and reported alongside CSA.

\section{CSA Theory Framework}
\label{sec:supp-theory}

\subsection{Matched Groups}

A matched group $G$ contains a positive arm $P_G$ and a negative arm $N_G$.  The arms are constructed so that they differ in the semantic factor that defines the task while remaining close in presentation and context.  For a group with $m$ matched contrasts,
\[
G=\{(x_j^+,x_j^-)\}_{j=1}^{m},\qquad
P_G=\{x_j^+\}_{j=1}^{m},\qquad
N_G=\{x_j^-\}_{j=1}^{m}.
\]
For SAT, $x_j^+$ is satisfiable and $x_j^-$ is a minimally edited unsatisfiable counterpart.  For vulnerability detection, $x_j^+$ is vulnerable and $x_j^-$ is a repaired or corresponding non-vulnerable counterpart.  In the SAT implementation used in this paper, a fixed problem size $N$ contains multiple groups $G\in\mathcal{G}_N$; the reported per-$N$ rows are obtained by first computing group-level quantities and then averaging them over $\mathcal{G}_N$.

\subsection{Arm Sensitivities and Factorized Separability}

Define the true positive and negative recalls
\[
\theta_P(G)=\Pr[f(X)=1\mid X\in P_G],
\qquad
\theta_N(G)=\Pr[f(X)=0\mid X\in N_G].
\]
The target separability ability is
\[
s_G=\theta_P(G)\theta_N(G).
\]
The empirical estimates are
\[
\hat\theta_P(G)=
\frac{1}{|P_G|}
\sum_{x\in P_G}\mathbf{1}\{f(x)=1\},
\qquad
\hat\theta_N(G)=
\frac{1}{|N_G|}
\sum_{x\in N_G}\mathbf{1}\{f(x)=0\},
\]
and
\[
\hat{s}_G=\hat\theta_P(G)\hat\theta_N(G).
\]

The product form is deliberate.  A one-sided all-positive predictor has high $\hat\theta_P$ and low $\hat\theta_N$; an all-negative predictor has the reverse.  In both cases, $\hat{s}_G$ is small.  Thus, factorized separability prevents single-arm success from being counted as strong matched reasoning evidence.

For a fixed problem size $N$, the reported arm summaries and separability point estimate are
\[
\hat\theta_P(N)=\frac{1}{|\mathcal{G}_N|}\sum_{G\in\mathcal{G}_N}\hat\theta_P(G),\qquad
\hat\theta_N(N)=\frac{1}{|\mathcal{G}_N|}\sum_{G\in\mathcal{G}_N}\hat\theta_N(G),
\]
\[
\hat s(N)=\frac{1}{|\mathcal{G}_N|}\sum_{G\in\mathcal{G}_N}\hat s_G
=\frac{1}{|\mathcal{G}_N|}\sum_{G\in\mathcal{G}_N}\hat\theta_P(G)\hat\theta_N(G).
\]
Therefore the per-$N$ separability point estimate is the average of within-group products, not the product of the averaged recalls; in general $\hat s(N)\neq \hat\theta_P(N)\hat\theta_N(N)$.

\subsection{Clopper--Pearson Lower Bound}

For $k$ successes in $n$ Bernoulli trials, the one-sided Clopper--Pearson lower bound at level $1-\gamma$ is
\[
L^{\mathrm{CP}}_{\gamma}(k,n)=
\begin{cases}
0, & k=0,\\
\mathrm{Beta}^{-1}(\gamma;k,n-k+1), & k>0,
\end{cases}
\]
where $\mathrm{Beta}^{-1}$ is the inverse beta cumulative distribution function.  In CSA, the confidence level is split across the two arms:
\[
L_P=L^{\mathrm{CP}}_{\alpha/2}(k_P,n_P),
\qquad
L_N=L^{\mathrm{CP}}_{\alpha/2}(k_N,n_N),
\]
where $k_P$ is the number of positive instances predicted positive and $k_N$ is the number of negative instances predicted negative.  The primary group-level estimator is
\[
\CSACP(G;\alpha)=L_PL_N.
\]
The reported per-$N$ CSA score is then the average over groups:
\[
\CSACP(N;\alpha)=\frac{1}{|\mathcal{G}_N|}\sum_{G\in\mathcal{G}_N}\CSACP(G;\alpha).
\]

\begin{proposition}[Coverage of the factorized CP lower bound]
Assume the within-group outcomes are Bernoulli draws satisfying the stated sampling model.  Then
\[
\Pr\left(s_G\ge \CSACP(G;\alpha)\right)\ge 1-\alpha.
\]
\end{proposition}

\begin{proof}
By construction,
\[
\Pr(\theta_P\ge L_P)\ge 1-\alpha/2,\qquad
\Pr(\theta_N\ge L_N)\ge 1-\alpha/2.
\]
By the union bound, both inequalities hold simultaneously with probability at least $1-\alpha$.  On that joint event,
\[
s_G=\theta_P\theta_N\ge L_PL_N=\CSACP(G;\alpha).
\]
\end{proof}

\subsection{Wilson Lower Bound}

The Wilson lower bound is reported as a robustness check.  For $\hat{p}=k/n$ and standard normal quantile $z_{1-\gamma}$, the one-sided Wilson lower bound is
\[
L^{\mathrm{W}}_{\gamma}(k,n)=
\frac{
\hat{p}+\frac{z^2}{2n}
-
z\sqrt{\frac{\hat{p}(1-\hat{p})}{n}+\frac{z^2}{4n^2}}
}{
1+\frac{z^2}{n}
}.
\]
CSA uses
\[
\CSAW(G;\alpha)=
L^{\mathrm{W}}_{\alpha/2}(k_P,n_P)
L^{\mathrm{W}}_{\alpha/2}(k_N,n_N).
\]
The Wilson variant is usually slightly less conservative than CP and helps diagnose whether conclusions are artifacts of a single interval construction.  The reported per-$N$ Wilson row is the corresponding group average
\[
\CSAW(N;\alpha)=\frac{1}{|\mathcal{G}_N|}\sum_{G\in\mathcal{G}_N}\CSAW(G;\alpha).
\]

\subsection{Formal Comparison with Existing Metrics}
\label{sec:supp-metric-rel}

We make precise the placement of $\hat s_G=\hat\theta_P\hat\theta_N$ among standard two-class metrics.  Throughout, $\theta_P$ denotes positive recall and $\theta_N$ denotes negative recall; in the standard binary-classification terminology, these correspond to sensitivity and specificity.

\paragraph{Order-equivalence to G-mean.}
Since $\Gmean=\sqrt{\theta_P\theta_N}=\sqrt{\hat s_G}$ and $t\mapsto\sqrt{t}$ is strictly increasing on $[0,1]$, $\hat s_G$ and $\Gmean$ induce the same total order on any set of classifiers.  Hence no ranking claim distinguishes them; the product scale is chosen only for the marginal interpretation, the clean per-arm confidence decomposition, and the $\tau=0.25$ reading discussed in the main text.

\paragraph{The zero-veto property within the power-mean family.}
Consider the symmetric power-mean aggregators
\[
M_p(\theta_P,\theta_N)=\left(\tfrac{1}{2}\theta_P^{\,p}+\tfrac{1}{2}\theta_N^{\,p}\right)^{1/p},\qquad p\neq 0,
\]
with $M_0$ defined as the limit $M_0=\sqrt{\theta_P\theta_N}=\Gmean$.  Say an aggregator has the \emph{zero-veto} property if it equals $0$ whenever either arm is $0$, i.e.\ a model that fails one arm entirely cannot be credited with two-sided ability.

\begin{proposition}[Power means and the veto property]\label{prop:veto}
$M_p$ has the zero-veto property if and only if $p\le 0$.  In particular the geometric mean ($p=0$) and the product $\hat s_G=M_0^2$ satisfy it, while the arithmetic mean $\BA=M_1$ does not.
\end{proposition}
\begin{proof}
For $p>0$ and $\theta_N=0$, $M_p(\theta_P,0)=(\theta_P^{\,p}/2)^{1/p}=2^{-1/p}\theta_P>0$ whenever $\theta_P>0$, so the veto fails (e.g.\ $\BA(\theta_P,0)=\theta_P/2$).  For $p\le 0$, $M_p$ is the limit/continuation that is dominated by the smaller argument; if either argument is $0$ then $M_p=0$.  The geometric mean is the $p\to0$ case and gives $M_0=\sqrt{\theta_P\theta_N}=0$ when either factor is $0$; squaring preserves this.
\end{proof}

Balanced accuracy and Youden's $J=2\BA-1$ are affine images of $M_1$ and therefore inherit its failure of the veto property: a one-sided predictor receives $\BA=0.5$, $J=0$ in the symmetric case, and more generally the same value as a balanced model of equal mean recall.  This is the formal sense in which additive summaries are unsuitable as a \emph{primary} two-sided certificate, and it is shared by CSA and G-mean.

\paragraph{Class-balance dependence of MCC.}
Let an evaluation set contain a fraction $\rho$ of positive instances with class-conditional recalls $(\theta_P,\theta_N)$.  The confusion-matrix entries are $\mathrm{TP}=\rho\theta_P$, $\mathrm{FN}=\rho(1-\theta_P)$, $\mathrm{TN}=(1-\rho)\theta_N$, $\mathrm{FP}=(1-\rho)(1-\theta_N)$ (per unit mass), and
\[
\MCC=\frac{\mathrm{TP}\cdot\mathrm{TN}-\mathrm{FP}\cdot\mathrm{FN}}
{\sqrt{(\mathrm{TP}{+}\mathrm{FP})(\mathrm{TP}{+}\mathrm{FN})(\mathrm{TN}{+}\mathrm{FP})(\mathrm{TN}{+}\mathrm{FN})}}.
\]

\begin{proposition}[MCC is not balance-invariant]\label{prop:mcc}
$\hat s_G$, $\Gmean$, $\BA$, and $J$ are functions of $(\theta_P,\theta_N)$ alone and hence invariant to $\rho$.  $\MCC$ is not: for $(\theta_P,\theta_N)=(0.8,0.6)$, $\MCC=0.408$ at $\rho=0.5$ and $\MCC=0.378$ at $\rho=0.75$.
\end{proposition}
\begin{proof}
Invariance of the recall-based metrics is immediate since their formulas contain no $\rho$.  For MCC, substitute the entries above.  At $\rho=0.5$ (e.g.\ $100$ positives, $100$ negatives): $\mathrm{TP}{=}80,\mathrm{FN}{=}20,\mathrm{TN}{=}60,\mathrm{FP}{=}40$, giving numerator $80\cdot60-40\cdot20=4000$ and denominator $\sqrt{120\cdot100\cdot100\cdot80}=9797.96$, so $\MCC=0.408$.  At $\rho=0.75$ ($300$ positives, $100$ negatives): $\mathrm{TP}{=}240,\mathrm{FN}{=}60,\mathrm{TN}{=}60,\mathrm{FP}{=}40$, giving numerator $240\cdot60-40\cdot60=12000$ and denominator $\sqrt{280\cdot300\cdot100\cdot120}=31749.0$, so $\MCC=0.378$.
\end{proof}

Because the matched groups are balanced by construction and the certification target is a balance-invariant two-sided recall, this dependence is the reason MCC is reported, if at all, as a complementary summary rather than the primary statistic.

\subsection{Group-Level Cluster Bootstrap}
\label{sec:supp-bootstrap}

The per-group CP and Wilson bounds assume binomial sampling within each group.  Generated instances share seed families and a fixed prompt template, so the appropriate uncertainty unit is the matched pair, not the instance.  We use a nonparametric cluster bootstrap.

\begin{algorithmic}[1]
\Require Matched pairs $\{(x_j^+,x_j^-)\}_{j=1}^m$, model $f$, level $1-\alpha$, replicates $B$
\For{$b=1$ to $B$}
  \State Sample $m$ pair indices $j_1,\ldots,j_m$ with replacement
  \State $\hat\theta_P^{(b)}\gets$ positive-side accuracy on $\{x_{j_t}^+\}$; $\hat\theta_N^{(b)}\gets$ negative-side accuracy on $\{x_{j_t}^-\}$
  \State $\hat s_G^{(b)}\gets \hat\theta_P^{(b)}\hat\theta_N^{(b)}$
\EndFor
\State \Return group-bootstrap replicates $\{\hat s_G^{(b)}\}_{b=1}^{B}$
\end{algorithmic}

For a fixed size $N$, the bootstrap first resamples matched pairs within each group $G\in\mathcal{G}_N$, computes $\hat s_G^{(b)}$ for every group, and then averages them:
\[
\hat s^{(b)}(N)=\frac{1}{|\mathcal{G}_N|}\sum_{G\in\mathcal{G}_N}\hat s_G^{(b)}.
\]
We then set $\CSAboot(N)$ to the empirical $\alpha$-quantile of $\{\hat s^{(b)}(N)\}_{b=1}^{B}$.  Resampling pairs (rather than instances) preserves both within-pair correlation and the shared-template structure within a group, so $\CSAboot$ is robust to the dependence that pooled instance-level intervals ignore.  We recommend $B\ge 2000$ and the bias-corrected and accelerated (BCa) quantile when the bootstrap distribution is skewed near the boundaries $0$ or $1$.  In reporting, $\CSAboot$ is placed beside $\CSACP$: when the two agree, the per-group binomial approximation is adequate for that configuration; when $\CSAboot<\CSACP$, the interval aggregation is optimistic and $\CSAboot$ should be used for the boundary decision $N^*(\tau)$.

\subsection{Threshold-Parametric Boundaries}

CSA is a continuous separability score in $[0,1]$.  To report capability boundaries across problem sizes, we define
\[
N^*(\tau)=\max\{N:\CSACP(N)\ge \tau\}.
\]
In the main paper, $\tau=0.25$ is used as the default moderate two-sided threshold.  In the symmetric case, $\theta_P=\theta_N=0.5$ gives $s_G=0.25$.  This threshold is not a baseline subtraction and should not be treated as a universal scientific constant.  It is an interpretable reporting threshold.  The full $\CSACP$ values are reported so readers can choose $\tau$ based on application risk.

For monotonic boundary summaries, we use a prefix convention: a model is said to certify ``through $N$'' at threshold $\tau$ when all reported sizes from the smallest tested size through $N$ satisfy $\CSACP\ge\tau$.  This prevents isolated non-monotone recoveries at larger sizes from being interpreted as sustained capability.

\section{Computation Protocol}
\label{sec:supp-protocol}

\subsection{SAT}

For each model and problem size, the SAT tables report:
\[
\mathrm{Acc},\quad
\mathrm{F1}_{\mathrm{SAT}},\quad
\mathrm{F1}_{\mathrm{UNSAT}},\quad
\mathrm{ADR},\quad
\mathrm{ADR}^{+w},\quad
\hat{s},\quad
\CSACP,\quad
\CSAW.
\]
$\mathrm{ADR}^{+w}$ is a SAT-side witness-aware reference: a SAT prediction counts only if the emitted satisfying assignment satisfies the formula.  CSA itself does not use witness information.  This separation is intentional.  It lets the paper compare constructive SAT-side evidence with category-only two-sided separability.

\subsection{Vulnerability Detection}

For vulnerability detection, the positive arm contains vulnerable functions and the negative arm contains non-vulnerable or repaired counterparts.  The same factorized global estimate is used:
\[
\hat{s}_{\mathrm{VD}}=\hat\theta_P\hat\theta_N,
\qquad
\CSA^{\mathrm{CP}}_{\mathrm{VD}}=
L^{\mathrm{CP}}_{\alpha/2}(k_P,n_P)
L^{\mathrm{CP}}_{\alpha/2}(k_N,n_N).
\]
The VD setting is useful because it removes the possibility that the method is SAT-specific.  It also illustrates the main failure mode of one-sided training: high vulnerable recall can coexist with very low non-vulnerable recognition, making the product lower bound small.

\section{Complete CSA Results}
\label{sec:supp-results}

This section provides the exact numerical backing for the main-paper CSA summaries.  The main paper visualizes the commercial 3-SAT trajectories and the open-weight 2-SAT trajectories as full-width figures; \Cref{tab:com_pern_csa,tab:os_pern_csa} list every plotted point, including standard descriptive metrics and the confidence-bounded separability metrics.  The boundary tables below summarize those same values under several threshold choices.

\subsection{Threshold Boundaries}

\begin{table*}[ht]
\centering
\caption{Prefix boundaries for commercial 3-SAT models under several CSA thresholds.  The main paper uses $\tau=0.25$ as the default moderate two-sided threshold.}
\label{tab:supp-boundary-3sat}
\renewcommand{\arraystretch}{1.12}
\resizebox{\textwidth}{!}{%
\begin{tabular}{lrrrr}
\toprule
Model & $N^*(0.50)$ & $N^*(0.25)$ & $N^*(0.10)$ & $N^*(0.05)$ \\
\midrule
GPT-5 & 25 & 30 & 40 & 60 \\
Opus 4.7 medium-thinking & 25 & 25 & 50 & 60 \\
DeepSeek-Reasoner & 10 & 10 & 25 & 25 \\
\bottomrule
\end{tabular}%
}
\end{table*}

\begin{table*}[ht]
\centering
\caption{Prefix boundaries for open-weight thinking-mode 2-SAT models under several CSA thresholds.}
\label{tab:supp-boundary-2sat}
\renewcommand{\arraystretch}{1.12}
\resizebox{\textwidth}{!}{%
\begin{tabular}{lrrrr}
\toprule
Model & $N^*(0.50)$ & $N^*(0.25)$ & $N^*(0.10)$ & $N^*(0.05)$ \\
\midrule
Qwen3-14B thinking & 15 & 20 & 50 & 50 \\
Qwen3-32B thinking & 10 & 15 & 20 & 30 \\
QwQ-32B thinking & 10 & 15 & 20 & 25 \\
\bottomrule
\end{tabular}%
}
\end{table*}

\subsection{Commercial 3-SAT CSA}

\begin{table*}[ht]
\centering
\caption{Per-$N$ confidence-bounded separability ability for the three commercial reasoning models on 3-SAT.  This table provides the exact values plotted in the main-paper commercial 3-SAT CSA figure.  Model abbreviations are: GPT-5; Opus 4.7 (med.\ think) = Claude Opus 4.7 medium-thinking; DeepSeek-Reasoner = DeepSeek-Reasoner.  $R_{\mathrm{SAT}}$, $R_{\mathrm{UNSAT}}$, and direct separability ability are direct pooled quantities computed by ignoring the internal group partition at each fixed $N$, where direct separability ability is their product.  By contrast, $\hat{s}$ is the group-averaged separability point estimate used by the Case-D analysis.  G-mean, balanced accuracy (BA), Youden's $J$, and MCC are companion descriptive metrics; $\CSACP$ is the primary Clopper--Pearson factorized lower bound, and $\CSAW$ is the Wilson robustness check.  The last three columns append the new pair-scenario results: $\mathrm{ADR}_{\mathrm{pair}}$ for the original fixed pairing, $\mathrm{ADR}_{\max}$ for the maximum rematching branch, and $\mathrm{ADR}_{\min}$ for the minimum rematching branch.}
\label{tab:com_pern_csa}
\renewcommand{\arraystretch}{1.12}
\resizebox{\textwidth}{!}{%
\begin{tabular}{p{1.7cm}r rrr rr r rr rrrrrrr rrr}
\toprule
Model & $N$ & Acc & F1$_{\mathrm{SAT}}$ & F1$_{\mathrm{UNSAT}}$ & ADR & ADR$^{+w}$ & $\hat{s}$ & $\CSACP$ & $\CSAW$ & $R_{\mathrm{SAT}}$ & $R_{\mathrm{UNSAT}}$ & $s_{\mathrm{direct}}$ & G-mean & BA & $J$ & MCC & $\mathrm{ADR}_{\mathrm{pair}}$ & $\mathrm{ADR}_{\max}$ & $\mathrm{ADR}_{\min}$ \\
\midrule
\multirow{7}{*}{\makecell[l]{\textbf{GPT-5}}}
  & 5  & 1.000 & 1.000 & 1.000 & 1.000 & 1.000 & 1.000 & \cellcolor{green!15}0.929 & \cellcolor{green!15}0.927 & 1.000 & 1.000 & 1.000 & 1.000 & 1.000 & 1.000 & 1.000 & 1.000 & 1.000 & 1.000 \\
  & 10 & 0.990 & 0.991 & 0.989 & 0.989 & 0.989 & 0.981 & \cellcolor{green!15}0.897 & \cellcolor{green!15}0.897 & 0.982 & 1.000 & 0.982 & 0.991 & 0.991 & 0.982 & 0.980 & 0.989 & 0.989 & 0.989 \\
  & 25 & 0.864 & 0.847 & 0.877 & 0.736 & 0.673 & 0.732 & \cellcolor{green!15}0.612 & \cellcolor{green!15}0.615 & 0.755 & 0.973 & 0.734 & 0.857 & 0.864 & 0.727 & 0.745 & 0.736 & 0.755 & 0.727 \\
  & 30 & 0.735 & 0.651 & 0.787 & 0.486 & 0.440 & 0.480 & \cellcolor{green!15}0.367 & \cellcolor{green!15}0.372 & 0.495 & 0.973 & 0.482 & 0.694 & 0.734 & 0.468 & 0.533 & 0.486 & 0.495 & 0.468 \\
  & 40 & 0.645 & 0.466 & 0.735 & 0.309 & 0.164 & 0.303 & \cellcolor{green!15}0.210 & \cellcolor{green!15}0.216 & 0.309 & 0.982 & 0.303 & 0.551 & 0.645 & 0.291 & 0.393 & 0.309 & 0.309 & 0.291 \\
  & 50 & 0.573 & 0.288 & 0.695 & 0.173 & 0.064 & 0.168 & \cellcolor{green!15}0.099 & \cellcolor{green!15}0.105 & 0.173 & 0.973 & 0.168 & 0.410 & 0.573 & 0.145 & 0.242 & 0.173 & 0.173 & 0.145 \\
  & 60 & 0.555 & 0.222 & 0.688 & 0.127 & 0.000 & 0.125 & \cellcolor{green!15}0.067 & \cellcolor{green!15}0.072 & 0.127 & 0.982 & 0.125 & 0.353 & 0.555 & 0.109 & 0.210 & 0.127 & 0.127 & 0.118 \\
\midrule
\multirow{5}{*}{\makecell[l]{\textbf{Opus 4.7}\\\textbf{(med.\ think)}}}
  & 5  & 0.945 & 0.947 & 0.942 & 0.897 & 0.813 & 0.893 & \cellcolor{green!15}0.782 & \cellcolor{green!15}0.784 & 0.982 & 0.907 & 0.891 & 0.944 & 0.945 & 0.889 & 0.892 & 0.897 & 0.907 & 0.888 \\
  & 10 & 0.959 & 0.960 & 0.958 & 0.917 & 0.861 & 0.919 & \cellcolor{green!15}0.816 & \cellcolor{green!15}0.817 & 0.982 & 0.936 & 0.919 & 0.958 & 0.959 & 0.917 & 0.918 & 0.917 & 0.926 & 0.917 \\
  & 25 & 0.935 & 0.956 & 0.875 & 0.759 & 0.655 & 0.783 & \cellcolor{green!15}0.592 & \cellcolor{green!15}0.602 & 0.989 & 0.800 & 0.791 & 0.889 & 0.894 & 0.789 & 0.838 & 0.759 & 0.759 & 0.759 \\
  & 50 & 0.527 & 0.620 & 0.372 & 0.200 & 0.060 & 0.230 & \cellcolor{green!15}0.124 & \cellcolor{green!15}0.131 & 0.853 & 0.256 & 0.218 & 0.467 & 0.555 & 0.109 & 0.134 & 0.200 & 0.220 & 0.140 \\
  & 60 & 0.530 & 0.656 & 0.256 & 0.156 & 0.000 & 0.152 & \cellcolor{green!15}0.078 & \cellcolor{green!15}0.084 & 0.943 & 0.155 & 0.146 & 0.382 & 0.549 & 0.098 & 0.157 & 0.156 & 0.169 & 0.117 \\
\midrule
\multirow{7}{*}{\makecell[l]{\textbf{DeepSeek-}\\\textbf{Reasoner}}}
  & 5  & 0.986 & 0.986 & 0.987 & 0.973 & 0.973 & 0.973 & \cellcolor{green!15}0.892 & \cellcolor{green!15}0.892 & 0.973 & 1.000 & 0.973 & 0.986 & 0.986 & 0.973 & 0.973 & 0.973 & 0.973 & 0.973 \\
  & 10 & 0.876 & 0.865 & 0.886 & 0.780 & 0.780 & 0.761 & \cellcolor{green!15}0.641 & \cellcolor{green!15}0.644 & 0.769 & 0.990 & 0.761 & 0.872 & 0.879 & 0.759 & 0.774 & 0.780 & 0.780 & 0.770 \\
  & 25 & 0.626 & 0.454 & 0.716 & 0.289 & 0.144 & 0.290 & \cellcolor{green!15}0.197 & \cellcolor{green!15}0.203 & 0.305 & 0.960 & 0.293 & 0.541 & 0.633 & 0.265 & 0.349 & 0.289 & 0.299 & 0.258 \\
  & 30 & 0.515 & 0.154 & 0.660 & 0.073 & 0.010 & 0.082 & \cellcolor{green!15}0.036 & \cellcolor{green!15}0.041 & 0.087 & 0.950 & 0.083 & 0.288 & 0.519 & 0.038 & 0.075 & 0.073 & 0.073 & 0.052 \\
  & 40 & 0.524 & 0.183 & 0.664 & 0.093 & 0.031 & 0.104 & \cellcolor{green!15}0.048 & \cellcolor{green!15}0.053 & 0.108 & 0.933 & 0.101 & 0.317 & 0.520 & 0.041 & 0.072 & 0.093 & 0.103 & 0.082 \\
  & 50 & 0.486 & 0.194 & 0.622 & 0.090 & 0.000 & 0.102 & \cellcolor{green!15}0.052 & \cellcolor{green!15}0.057 & 0.120 & 0.873 & 0.105 & 0.324 & 0.496 & -0.007 & -0.011 & 0.090 & 0.110 & 0.030 \\
  & 60 & 0.509 & 0.160 & 0.653 & 0.086 & 0.010 & 0.086 & \cellcolor{green!15}0.039 & \cellcolor{green!15}0.044 & 0.093 & 0.934 & 0.086 & 0.294 & 0.513 & 0.027 & 0.049 & 0.086 & 0.086 & 0.038 \\
\bottomrule
\end{tabular}%
}
\end{table*}

The commercial 3-SAT table now exposes the descriptive structure more transparently.  The added Recall$_{\mathrm{SAT}}$ and Recall$_{\mathrm{UNSAT}}$ columns show exactly which side drives each change, while the direct separability ability column reports their pooled product when the within-$N$ group partition is ignored.  By construction, this direct separability ability is the square of G-mean, and the reported values follow that identity up to rounding.  Even with those additional descriptive views, the main inferential point is unchanged: the CSA lower bounds decline faster because they certify two-sided separability under finite-sample conservatism rather than merely summarizing pooled correctness.  MCC remains useful as a complementary summary because it penalizes degenerate or strongly imbalanced confusion patterns more sharply than accuracy alone.  GPT-5 and Opus 4.7 retain moderate direct descriptive values at larger sizes while their certified lower bounds are already much smaller, whereas DeepSeek-Reasoner shows the same collapse earlier through its rapidly weakening Recall$_{\mathrm{SAT}}$ and near-zero MCC.

\subsection{Open-Weight 2-SAT CSA}

\begin{table*}[ht]
\centering
\caption{Per-$N$ confidence-bounded separability ability for the six open-weight thinking/reasoning-mode models on 2-SAT.  This table provides the exact values plotted in the main-paper open-weight 2-SAT CSA figure.  Model abbreviations are: gpt-oss-20b (reason) = OpenAI gpt-oss-20b reasoning mode; Qwen3-14B/32B and QwQ-32B are the corresponding Qwen thinking-mode runs; Llama-Nemotron-Super-49B (reason) = NVIDIA Llama-3.3-Nemotron-Super-49B-v1.5 reasoning run; Nemotron-H-47B-Reasoning (reason) = NVIDIA Nemotron-H-47B-Reasoning-128K reasoning run.  $R_{\mathrm{SAT}}$, $R_{\mathrm{UNSAT}}$, and direct separability ability are direct pooled quantities computed by ignoring the internal group partition at each fixed $N$, where direct separability ability is their product.  By contrast, $\hat{s}$ is the group-averaged separability point estimate used by the Case-D analysis.  G-mean, balanced accuracy (BA), Youden's $J$, and MCC are companion descriptive metrics.  $\ADRwu$ is a stricter witness-aware variant that requires both a verified SAT assignment and, on the UNSAT side, a verified UNSAT clause subset from the follow-up prompt.  `--' indicates that no eligible pair exists under this stricter denominator, that is, there is no pair at that $N$ with both a SAT-side assignment and an UNSAT-side clause-subset follow-up result.  The last three columns append the new pair-scenario results: $\mathrm{ADR}_{\mathrm{pair}}$ for the original fixed pairing, $\mathrm{ADR}_{\max}$ for the maximum rematching branch, and $\mathrm{ADR}_{\min}$ for the minimum rematching branch.}
\label{tab:os_pern_csa}
\renewcommand{\arraystretch}{1.12}
\resizebox{\textwidth}{!}{%
\begin{tabular}{p{1.75cm}r rrr rrr r rr rrrrrrr rrr}
\toprule
Model & $N$ & Acc & F1$_{\mathrm{SAT}}$ & F1$_{\mathrm{UNSAT}}$ & ADR & ADR$^{+w}$ & $\ADRwu$ & $\hat{s}$ & $\CSACP$ & $\CSAW$ & $R_{\mathrm{SAT}}$ & $R_{\mathrm{UNSAT}}$ & $s_{\mathrm{direct}}$ & G-mean & BA & $J$ & MCC & $\mathrm{ADR}_{\mathrm{pair}}$ & $\mathrm{ADR}_{\max}$ & $\mathrm{ADR}_{\min}$ \\
\midrule
\multirow{10}{*}{\makecell[l]{\textbf{gpt-oss-20b}\\\textbf{(reason)}}}
  & 3  & 1.000 & 1.000 & 1.000 & 1.000 & 1.000 & 0.818 & 1.000 & \cellcolor{green!15}0.976 & \cellcolor{green!15}0.969 & 1.000 & 1.000 & 1.000 & 1.000 & 1.000 & 1.000 & 1.000 & 1.000 & 1.000 & 1.000 \\
  & 5  & 1.000 & 1.000 & 1.000 & 1.000 & 1.000 & 0.818 & 1.000 & \cellcolor{green!15}0.976 & \cellcolor{green!15}0.969 & 1.000 & 1.000 & 1.000 & 1.000 & 1.000 & 1.000 & 1.000 & 1.000 & 1.000 & 1.000 \\
  & 8  & 0.955 & 0.952 & 0.957 & 0.909 & 0.864 & 0.444 & 0.909 & \cellcolor{green!15}0.863 & \cellcolor{green!15}0.855 & 0.909 & 1.000 & 0.909 & 0.953 & 0.955 & 0.909 & 0.913 & 0.909 & 0.909 & 0.909 \\
  & 10 & 0.841 & 0.844 & 0.837 & 0.727 & 0.682 & 0.333 & 0.707 & \cellcolor{green!15}0.631 & \cellcolor{green!15}0.620 & 0.864 & 0.818 & 0.707 & 0.841 & 0.841 & 0.682 & 0.683 & 0.727 & 0.818 & 0.682 \\
  & 15 & 0.773 & 0.737 & 0.800 & 0.545 & 0.500 & 0.000 & 0.562 & \cellcolor{green!15}0.485 & \cellcolor{green!15}0.475 & 0.636 & 0.909 & 0.579 & 0.761 & 0.773 & 0.545 & 0.567 & 0.545 & 0.636 & 0.545 \\
  & 20 & 0.773 & 0.706 & 0.815 & 0.545 & 0.409 & 0.000 & 0.545 & \cellcolor{green!15}0.467 & \cellcolor{green!15}0.457 & 0.545 & 1.000 & 0.545 & 0.739 & 0.773 & 0.545 & 0.612 & 0.545 & 0.545 & 0.545 \\
  & 25 & 0.568 & 0.457 & 0.642 & 0.227 & 0.182 & --    & 0.281 & \cellcolor{green!15}0.216 & \cellcolor{green!15}0.210 & 0.364 & 0.773 & 0.281 & 0.530 & 0.568 & 0.136 & 0.149 & 0.227 & 0.364 & 0.182 \\
  & 30 & 0.568 & 0.296 & 0.689 & 0.136 & 0.045 & --    & 0.169 & \cellcolor{green!15}0.118 & \cellcolor{green!15}0.116 & 0.182 & 0.955 & 0.174 & 0.417 & 0.568 & 0.136 & 0.215 & 0.136 & 0.182 & 0.136 \\
  & 40 & 0.545 & 0.286 & 0.667 & 0.182 & 0.000 & 0.000 & 0.165 & \cellcolor{green!15}0.112 & \cellcolor{green!15}0.110 & 0.182 & 0.909 & 0.165 & 0.407 & 0.545 & 0.091 & 0.132 & 0.182 & 0.182 & 0.091 \\
  & 50 & 0.500 & 0.083 & 0.656 & 0.045 & 0.000 & --    & 0.045 & \cellcolor{green!15}0.026 & \cellcolor{green!15}0.026 & 0.045 & 0.955 & 0.043 & 0.208 & 0.500 & 0.000 & 0.000 & 0.045 & 0.045 & 0.045 \\
\midrule
\multirow{10}{*}{\makecell[l]{\textbf{Qwen3-14B}\\\textbf{(think)}}}
  & 3  & 0.990 & 0.990 & 0.990 & 0.981 & 0.981 & 0.118 & 0.981 & \cellcolor{green!15}0.950 & \cellcolor{green!15}0.943 & 0.981 & 1.000 & 0.981 & 0.990 & 0.990 & 0.981 & 0.981 & 0.981 & 0.981 & 0.981 \\
  & 5  & 0.974 & 0.973 & 0.975 & 0.948 & 0.942 & 0.142 & 0.948 & \cellcolor{green!15}0.909 & \cellcolor{green!15}0.901 & 0.948 & 1.000 & 0.948 & 0.974 & 0.974 & 0.948 & 0.949 & 0.948 & 0.948 & 0.948 \\
  & 8  & 0.974 & 0.973 & 0.975 & 0.948 & 0.942 & 0.179 & 0.948 & \cellcolor{green!15}0.907 & \cellcolor{green!15}0.899 & 0.948 & 1.000 & 0.948 & 0.974 & 0.974 & 0.948 & 0.949 & 0.948 & 0.948 & 0.948 \\
  & 10 & 0.942 & 0.938 & 0.945 & 0.883 & 0.877 & 0.112 & 0.883 & \cellcolor{green!15}0.833 & \cellcolor{green!15}0.824 & 0.883 & 1.000 & 0.883 & 0.940 & 0.942 & 0.883 & 0.889 & 0.883 & 0.883 & 0.883 \\
  & 15 & 0.844 & 0.815 & 0.865 & 0.688 & 0.656 & 0.085 & 0.688 & \cellcolor{green!15}0.620 & \cellcolor{green!15}0.611 & 0.688 & 1.000 & 0.688 & 0.830 & 0.844 & 0.688 & 0.724 & 0.688 & 0.688 & 0.688 \\
  & 20 & 0.669 & 0.519 & 0.748 & 0.338 & 0.286 & 0.061 & 0.348 & \cellcolor{green!15}0.282 & \cellcolor{green!15}0.276 & 0.357 & 0.981 & 0.350 & 0.592 & 0.669 & 0.338 & 0.432 & 0.338 & 0.357 & 0.338 \\
  & 25 & 0.604 & 0.371 & 0.711 & 0.214 & 0.156 & 0.000 & 0.230 & \cellcolor{green!15}0.179 & \cellcolor{green!15}0.175 & 0.234 & 0.974 & 0.228 & 0.477 & 0.604 & 0.208 & 0.309 & 0.214 & 0.234 & 0.208 \\
  & 30 & 0.516 & 0.280 & 0.636 & 0.162 & 0.058 & 0.000 & 0.153 & \cellcolor{green!15}0.105 & \cellcolor{green!15}0.103 & 0.188 & 0.844 & 0.159 & 0.399 & 0.516 & 0.032 & 0.043 & 0.162 & 0.188 & 0.078 \\
  & 40 & 0.490 & 0.338 & 0.585 & 0.203 & 0.000 & 0.000 & 0.184 & \cellcolor{green!15}0.136 & \cellcolor{green!15}0.133 & 0.255 & 0.734 & 0.187 & 0.433 & 0.495 & -0.011 & -0.012 & 0.203 & 0.238 & 0.077 \\
  & 50 & 0.479 & 0.313 & 0.580 & 0.154 & 0.000 & 0.000 & 0.163 & \cellcolor{green!15}0.118 & \cellcolor{green!15}0.115 & 0.238 & 0.720 & 0.171 & 0.414 & 0.479 & -0.042 & -0.048 & 0.154 & 0.231 & 0.063 \\
\midrule
\multirow{10}{*}{\makecell[l]{\textbf{Qwen3-32B}\\\textbf{(think)}}}
  & 3  & 0.979 & 0.979 & 0.980 & 0.959 & 0.959 & 0.222 & 0.959 & \cellcolor{green!15}0.922 & \cellcolor{green!15}0.914 & 0.959 & 1.000 & 0.959 & 0.979 & 0.979 & 0.959 & 0.959 & 0.959 & 0.959 & 0.959 \\
  & 5  & 0.955 & 0.952 & 0.957 & 0.909 & 0.909 & 0.355 & 0.909 & \cellcolor{green!15}0.864 & \cellcolor{green!15}0.855 & 0.909 & 1.000 & 0.909 & 0.953 & 0.955 & 0.909 & 0.913 & 0.909 & 0.909 & 0.909 \\
  & 8  & 0.926 & 0.920 & 0.931 & 0.851 & 0.843 & 0.100 & 0.851 & \cellcolor{green!15}0.798 & \cellcolor{green!15}0.790 & 0.851 & 1.000 & 0.851 & 0.923 & 0.926 & 0.851 & 0.861 & 0.851 & 0.851 & 0.851 \\
  & 10 & 0.905 & 0.895 & 0.913 & 0.810 & 0.777 & 0.281 & 0.810 & \cellcolor{green!15}0.753 & \cellcolor{green!15}0.744 & 0.810 & 1.000 & 0.810 & 0.900 & 0.905 & 0.810 & 0.825 & 0.810 & 0.810 & 0.810 \\
  & 15 & 0.748 & 0.663 & 0.799 & 0.496 & 0.430 & 0.071 & 0.496 & \cellcolor{green!15}0.422 & \cellcolor{green!15}0.413 & 0.496 & 1.000 & 0.496 & 0.704 & 0.748 & 0.496 & 0.574 & 0.496 & 0.496 & 0.496 \\
  & 20 & 0.624 & 0.397 & 0.727 & 0.248 & 0.207 & 0.077 & 0.248 & \cellcolor{green!15}0.189 & \cellcolor{green!15}0.185 & 0.248 & 1.000 & 0.248 & 0.498 & 0.624 & 0.248 & 0.376 & 0.248 & 0.248 & 0.248 \\
  & 25 & 0.568 & 0.246 & 0.698 & 0.142 & 0.092 & 0.000 & 0.140 & \cellcolor{green!15}0.098 & \cellcolor{green!15}0.096 & 0.142 & 0.992 & 0.140 & 0.375 & 0.567 & 0.133 & 0.254 & 0.142 & 0.142 & 0.133 \\
  & 30 & 0.574 & 0.270 & 0.700 & 0.157 & 0.083 & 0.125 & 0.156 & \cellcolor{green!15}0.112 & \cellcolor{green!15}0.110 & 0.157 & 0.992 & 0.156 & 0.395 & 0.574 & 0.149 & 0.270 & 0.157 & 0.157 & 0.149 \\
  & 40 & 0.525 & 0.173 & 0.667 & 0.099 & 0.000 & 0.000 & 0.096 & \cellcolor{green!15}0.067 & \cellcolor{green!15}0.066 & 0.099 & 0.950 & 0.094 & 0.307 & 0.525 & 0.050 & 0.094 & 0.099 & 0.099 & 0.083 \\
  & 50 & 0.467 & 0.273 & 0.579 & 0.143 & 0.000 & 0.000 & 0.129 & \cellcolor{green!15}0.089 & \cellcolor{green!15}0.087 & 0.202 & 0.727 & 0.147 & 0.383 & 0.464 & -0.071 & -0.083 & 0.143 & 0.185 & 0.042 \\
\midrule
\multirow{10}{*}{\makecell[l]{\textbf{QwQ-32B}\\\textbf{(think)}}}
  & 3  & 1.000 & 1.000 & 1.000 & 1.000 & 1.000 & 0.156 & 1.000 & \cellcolor{green!15}0.976 & \cellcolor{green!15}0.969 & 1.000 & 1.000 & 1.000 & 1.000 & 1.000 & 1.000 & 1.000 & 1.000 & 1.000 & 1.000 \\
  & 5  & 0.990 & 0.990 & 0.990 & 0.980 & 0.980 & 0.194 & 0.980 & \cellcolor{green!15}0.949 & \cellcolor{green!15}0.942 & 0.980 & 1.000 & 0.980 & 0.990 & 0.990 & 0.980 & 0.980 & 0.980 & 0.980 & 0.980 \\
  & 8  & 0.980 & 0.980 & 0.980 & 0.960 & 0.960 & 0.111 & 0.960 & \cellcolor{green!15}0.924 & \cellcolor{green!15}0.916 & 0.970 & 0.990 & 0.960 & 0.980 & 0.980 & 0.960 & 0.960 & 0.960 & 0.960 & 0.960 \\
  & 10 & 0.949 & 0.947 & 0.952 & 0.899 & 0.869 & 0.071 & 0.899 & \cellcolor{green!15}0.857 & \cellcolor{green!15}0.849 & 0.899 & 1.000 & 0.899 & 0.948 & 0.949 & 0.899 & 0.904 & 0.899 & 0.899 & 0.899 \\
  & 15 & 0.798 & 0.747 & 0.832 & 0.596 & 0.545 & 0.087 & 0.596 & \cellcolor{green!15}0.521 & \cellcolor{green!15}0.511 & 0.596 & 1.000 & 0.596 & 0.772 & 0.798 & 0.596 & 0.652 & 0.596 & 0.596 & 0.596 \\
  & 20 & 0.641 & 0.441 & 0.736 & 0.283 & 0.242 & 0.231 & 0.283 & \cellcolor{green!15}0.222 & \cellcolor{green!15}0.218 & 0.283 & 1.000 & 0.283 & 0.532 & 0.641 & 0.283 & 0.406 & 0.283 & 0.283 & 0.283 \\
  & 25 & 0.576 & 0.276 & 0.700 & 0.162 & 0.111 & 0.000 & 0.161 & \cellcolor{green!15}0.122 & \cellcolor{green!15}0.119 & 0.162 & 0.990 & 0.160 & 0.400 & 0.576 & 0.152 & 0.270 & 0.162 & 0.162 & 0.152 \\
  & 30 & 0.515 & 0.059 & 0.673 & 0.030 & 0.000 & --    & 0.030 & \cellcolor{green!15}0.017 & \cellcolor{green!15}0.017 & 0.030 & 1.000 & 0.030 & 0.174 & 0.515 & 0.030 & 0.124 & 0.030 & 0.030 & 0.030 \\
  & 40 & 0.487 & 0.093 & 0.642 & 0.033 & 0.011 & 0.000 & 0.039 & \cellcolor{green!15}0.018 & \cellcolor{green!15}0.019 & 0.051 & 0.967 & 0.049 & 0.221 & 0.509 & 0.017 & 0.043 & 0.033 & 0.044 & 0.022 \\
  & 50 & 0.506 & 0.044 & 0.667 & 0.023 & 0.000 & --    & 0.023 & \cellcolor{green!15}0.013 & \cellcolor{green!15}0.013 & 0.023 & 0.989 & 0.022 & 0.150 & 0.506 & 0.011 & 0.044 & 0.023 & 0.023 & 0.023 \\
\midrule
\multirow{8}{*}{\makecell[l]{\textbf{Llama-Nemotron-}\\\textbf{Super-49B (reason)}}}
  & 3  & 1.000 & 1.000 & 1.000 & 1.000 & 0.909 & --    & 1.000 & \cellcolor{green!15}0.976 & \cellcolor{green!15}0.969 & 1.000 & 1.000 & 1.000 & 1.000 & 1.000 & 1.000 & 1.000 & 1.000 & 1.000 & 1.000 \\
  & 5  & 1.000 & 1.000 & 1.000 & 1.000 & 1.000 & --    & 1.000 & \cellcolor{green!15}0.976 & \cellcolor{green!15}0.969 & 1.000 & 1.000 & 1.000 & 1.000 & 1.000 & 1.000 & 1.000 & 1.000 & 1.000 & 1.000 \\
  & 8  & 0.864 & 0.842 & 0.880 & 0.727 & 0.727 & --    & 0.727 & \cellcolor{green!15}0.653 & \cellcolor{green!15}0.642 & 0.727 & 1.000 & 0.727 & 0.853 & 0.864 & 0.727 & 0.756 & 0.727 & 0.727 & 0.727 \\
  & 10 & 0.773 & 0.706 & 0.815 & 0.545 & 0.545 & --    & 0.545 & \cellcolor{green!15}0.467 & \cellcolor{green!15}0.457 & 0.545 & 1.000 & 0.545 & 0.739 & 0.773 & 0.545 & 0.612 & 0.545 & 0.545 & 0.545 \\
  & 15 & 0.727 & 0.625 & 0.786 & 0.455 & 0.364 & --    & 0.455 & \cellcolor{green!15}0.377 & \cellcolor{green!15}0.369 & 0.455 & 1.000 & 0.455 & 0.674 & 0.727 & 0.455 & 0.542 & 0.455 & 0.455 & 0.455 \\
  & 20 & 0.818 & 0.778 & 0.846 & 0.636 & 0.545 & --    & 0.636 & \cellcolor{green!15}0.558 & \cellcolor{green!15}0.548 & 0.636 & 1.000 & 0.636 & 0.798 & 0.818 & 0.636 & 0.683 & 0.636 & 0.636 & 0.636 \\
  & 25 & 0.773 & 0.706 & 0.815 & 0.545 & 0.182 & --    & 0.545 & \cellcolor{green!15}0.467 & \cellcolor{green!15}0.457 & 0.545 & 1.000 & 0.545 & 0.739 & 0.773 & 0.545 & 0.612 & 0.545 & 0.545 & 0.545 \\
  & 30 & 0.400 & 0.308 & 0.471 & 0.250 & 0.000 & --    & 0.182 & \cellcolor{green!15}0.094 & \cellcolor{green!15}0.095 & 0.182 & 1.000 & 0.182 & 0.426 & 0.591 & 0.182 & 0.237 & 0.250 & 0.250 & 0.250 \\
\midrule
\multirow{11}{*}{\makecell[l]{\textbf{Nemotron-H-}\\\textbf{47B-Reasoning}\\\textbf{(reason)}}}
  & 3  & 0.524 & 0.657 & 0.224 & 0.077 & 0.046 & --    & 0.080 & \cellcolor{green!15}0.069 & \cellcolor{green!15}0.068 & 0.911 & 0.137 & 0.125 & 0.354 & 0.524 & 0.048 & 0.076 & 0.077 & 0.087 & 0.074 \\
  & 5  & 0.510 & 0.668 & 0.066 & 0.023 & 0.004 & --    & 0.027 & \cellcolor{green!15}0.019 & \cellcolor{green!15}0.019 & 0.985 & 0.035 & 0.034 & 0.185 & 0.510 & 0.019 & 0.062 & 0.023 & 0.031 & 0.021 \\
  & 8  & 0.499 & 0.666 & --    & 0.000 & 0.000 & --    & 0.000 & \cellcolor{green!15}0.000 & \cellcolor{green!15}0.000 & 0.998 & 0.000 & 0.000 & 0.000 & 0.499 & -0.002 & -0.031 & 0.000 & 0.000 & 0.000 \\
  & 10 & 0.497 & 0.664 & 0.004 & 0.002 & 0.000 & --    & 0.002 & \cellcolor{green!15}0.001 & \cellcolor{green!15}0.001 & 0.992 & 0.002 & 0.002 & 0.044 & 0.497 & -0.006 & -0.042 & 0.002 & 0.002 & 0.002 \\
  & 15 & 0.513 & 0.668 & 0.087 & 0.027 & 0.000 & --    & 0.030 & \cellcolor{green!15}0.024 & \cellcolor{green!15}0.024 & 0.979 & 0.046 & 0.045 & 0.213 & 0.513 & 0.025 & 0.070 & 0.027 & 0.031 & 0.025 \\
  & 20 & 0.510 & 0.664 & 0.093 & 0.043 & 0.000 & --    & 0.043 & \cellcolor{green!15}0.037 & \cellcolor{green!15}0.036 & 0.969 & 0.050 & 0.049 & 0.221 & 0.510 & 0.019 & 0.049 & 0.043 & 0.046 & 0.041 \\
  & 25 & 0.526 & 0.666 & 0.183 & 0.070 & 0.000 & --    & 0.071 & \cellcolor{green!15}0.058 & \cellcolor{green!15}0.057 & 0.946 & 0.106 & 0.101 & 0.317 & 0.526 & 0.052 & 0.096 & 0.070 & 0.079 & 0.062 \\
  & 30 & 0.503 & 0.665 & 0.034 & 0.015 & 0.000 & --    & 0.017 & \cellcolor{green!15}0.012 & \cellcolor{green!15}0.012 & 0.988 & 0.017 & 0.017 & 0.131 & 0.503 & 0.006 & 0.024 & 0.015 & 0.017 & 0.015 \\
  & 40 & 0.527 & 0.673 & 0.147 & 0.074 & 0.000 & --    & 0.074 & \cellcolor{green!15}0.055 & \cellcolor{green!15}0.054 & 0.973 & 0.081 & 0.079 & 0.281 & 0.527 & 0.054 & 0.120 & 0.074 & 0.081 & 0.062 \\
  & 50 & 0.516 & 0.667 & 0.120 & 0.046 & 0.000 & --    & 0.052 & \cellcolor{green!15}0.042 & \cellcolor{green!15}0.041 & 0.967 & 0.066 & 0.064 & 0.252 & 0.516 & 0.033 & 0.076 & 0.046 & 0.062 & 0.041 \\
  & 60 & 0.537 & 0.674 & 0.206 & 0.105 & 0.000 & --    & 0.107 & \cellcolor{green!15}0.085 & \cellcolor{green!15}0.083 & 0.949 & 0.121 & 0.114 & 0.338 & 0.535 & 0.070 & 0.125 & 0.105 & 0.119 & 0.093 \\
\bottomrule
\end{tabular}%
}
\end{table*}

The open-weight table still shows the same broad pattern, but the refreshed 2-SAT rerun changes several details materially.  The added recall columns make clear that many early rows remain close to one-sided perfection on the UNSAT side, so the direct separability ability is still driven mainly by Recall$_{\mathrm{SAT}}$, with direct separability ability equal to the square of G-mean up to rounding.  The gpt-oss-20b block remains very strong at small $N$, but its stricter witness-aware columns are now nonzero rather than unavailable because corresponding follow-up outputs are present in the current workspace.  The Llama-Nemotron-Super-49B (reason) block is no longer limited to $N=3$ and $N=5$: under the refreshed NVIDIA reasoning-on results it now spans $N=3$ through $N=30$, with perfect two-sided rows at $N=3$ and $N=5$, then a gradual decline through larger $N$.  Nemotron-H-47B-Reasoning remains the weakest of the six models on matched two-sided separability, but the refreshed reasoning-on rerun no longer collapses to the earlier degenerate small-$N$ pattern in exactly the same way; instead, it shows very high SAT recall together with persistently weak UNSAT recall across all $N$, yielding small but generally nonzero direct products, CSA bounds, and pair-based ADR variants.  The $\ADRwu$ column is stricter still: it preserves a pair-level success only when the SAT assignment is verified and the UNSAT follow-up returns a verified unsatisfiable clause subset.  In the current workspace, nonzero $\ADRwu$ values appear for gpt-oss-20b, Qwen3-14B, Qwen3-32B, and QwQ-32B, while both NVIDIA blocks remain at $0$ throughout the reported rows.  MCC remains useful as a complementary view because it reacts strongly when one side of the confusion matrix dominates, helping distinguish balanced discrimination from superficially acceptable pooled accuracy.  Several rows still retain positive MCC and nonzero direct products while $\CSACP$ and $\CSAW$ are already much smaller, which remains the core point of the CSA analysis.

\subsection{Open-Weight 3-Coloring CSA}

\begin{table*}[ht]
\centering
\caption{Per-$N$ confidence-bounded separability ability for the six open-weight/open models on graph 3-coloring.  This auxiliary table applies the same matched-group CSA pipeline to a non-SAT decision problem.  Model abbreviations are: GPT-OSS-20B = OpenAI gpt-oss-20b thinking run; Llama-Nemo-49B = NVIDIA Llama-3.3-Nemotron-Super-49B-v1.5 reasoning-on run; Nemo-H-47B = NVIDIA Nemotron-H-47B-Reasoning-128K reasoning-on run; Qwen3-14B, Qwen3-32B, and QwQ-32B are the corresponding graph-coloring thinking-mode runs.  $R_{\mathrm{SAT}}$ and $R_{\mathrm{UNSAT}}$ denote recall on the 3-colorable and non-3-colorable sides respectively, but we keep the notation aligned with the rest of the paper for compactness.  Direct separability ability is the pooled within-$N$ product of those two recalls when the internal group partition is ignored.  By contrast, $\hat{s}$ is the group-averaged separability point estimate used by the Case-D analysis.  G-mean, balanced accuracy (BA), Youden's $J$, and MCC are companion descriptive metrics; $\CSACP$ is the primary Clopper--Pearson factorized lower bound, and $\CSAW$ is the Wilson robustness check.  The last three columns append the new pair-scenario results: $\mathrm{ADR}_{\mathrm{pair}}$ for the original fixed pairing, $\mathrm{ADR}_{\max}$ for the maximum rematching branch, and $\mathrm{ADR}_{\min}$ for the minimum rematching branch.}
\label{tab:gc_pern_csa}
\renewcommand{\arraystretch}{1.12}
\resizebox{\textwidth}{!}{%
\begin{tabular}{p{1.35cm}r rrr rr r rr rrrrrrr rrr}
\toprule
Model & $N$ & Acc & F1$_{\mathrm{SAT}}$ & F1$_{\mathrm{UNSAT}}$ & ADR & ADR$^{+w}$ & $\hat{s}$ & $\CSACP$ & $\CSAW$ & $R_{\mathrm{SAT}}$ & $R_{\mathrm{UNSAT}}$ & $s_{\mathrm{direct}}$ & G-mean & BA & $J$ & MCC & $\mathrm{ADR}_{\mathrm{pair}}$ & $\mathrm{ADR}_{\max}$ & $\mathrm{ADR}_{\min}$ \\
\midrule
\multirow{8}{*}{\makecell[l]{\textbf{GPT-OSS-20B}\\\textbf{(think)}}}
  & 8  & 0.953 & 0.952 & 0.955 & 0.905 & 0.857 & 0.909 & \cellcolor{green!15}0.852 & \cellcolor{green!15}0.843 & 0.909 & 1.000 & 0.909 & 0.953 & 0.955 & 0.909 & 0.911 & 0.905 & 0.905 & 0.905 \\
  & 10 & 0.886 & 0.878 & 0.894 & 0.818 & 0.818 & 0.781 & \cellcolor{green!15}0.710 & \cellcolor{green!15}0.700 & 0.818 & 0.955 & 0.781 & 0.884 & 0.886 & 0.773 & 0.780 & 0.818 & 0.818 & 0.773 \\
  & 15 & 0.545 & 0.286 & 0.667 & 0.182 & 0.136 & 0.165 & \cellcolor{green!15}0.130 & \cellcolor{green!15}0.127 & 0.182 & 0.909 & 0.165 & 0.407 & 0.545 & 0.091 & 0.132 & 0.182 & 0.182 & 0.136 \\
  & 20 & 0.659 & 0.462 & 0.750 & 0.316 & 0.316 & 0.313 & \cellcolor{green!15}0.242 & \cellcolor{green!15}0.236 & 0.300 & 1.000 & 0.300 & 0.548 & 0.650 & 0.300 & 0.424 & 0.322 & 0.322 & 0.322 \\
  & 30 & 0.528 & --    & 0.691 & 0.000 & 0.000 & 0.000 & \cellcolor{green!15}0.000 & \cellcolor{green!15}0.000 & 0.000 & 1.000 & 0.000 & 0.000 & 0.500 & 0.000 & --    & 0.000 & 0.000 & 0.000 \\
  & 40 & 0.571 & 0.118 & 0.717 & 0.071 & 0.000 & 0.062 & \cellcolor{green!15}0.036 & \cellcolor{green!15}0.036 & 0.062 & 1.000 & 0.062 & 0.250 & 0.531 & 0.062 & 0.187 & 0.062 & 0.062 & 0.062 \\
  & 50 & 0.440 & 0.222 & 0.562 & 0.182 & 0.000 & 0.149 & \cellcolor{green!15}0.098 & \cellcolor{green!15}0.096 & 0.143 & 0.818 & 0.117 & 0.342 & 0.481 & -0.039 & -0.053 & 0.182 & 0.182 & 0.000 \\
  & 60 & 0.550 & 0.308 & 0.667 & 0.111 & 0.000 & 0.180 & \cellcolor{green!15}0.120 & \cellcolor{green!15}0.117 & 0.200 & 0.900 & 0.180 & 0.424 & 0.550 & 0.100 & 0.140 & 0.111 & 0.222 & 0.111 \\
\midrule
\multirow{8}{*}{\makecell[l]{\textbf{Llama-Nemo-49B}\\\textbf{(r-on)}}}
  & 8  & 0.932 & 0.927 & 0.936 & 0.864 & 0.136 & 0.864 & \cellcolor{green!15}0.802 & \cellcolor{green!15}0.792 & 0.864 & 1.000 & 0.864 & 0.929 & 0.932 & 0.864 & 0.872 & 0.864 & 0.864 & 0.864 \\
  & 10 & 0.909 & 0.900 & 0.917 & 0.818 & 0.227 & 0.818 & \cellcolor{green!15}0.751 & \cellcolor{green!15}0.740 & 0.818 & 1.000 & 0.818 & 0.905 & 0.909 & 0.818 & 0.832 & 0.818 & 0.818 & 0.818 \\
  & 15 & 0.545 & 0.231 & 0.677 & 0.136 & 0.091 & 0.124 & \cellcolor{green!15}0.092 & \cellcolor{green!15}0.090 & 0.136 & 0.955 & 0.130 & 0.361 & 0.545 & 0.091 & 0.158 & 0.136 & 0.136 & 0.091 \\
  & 20 & 0.545 & 0.167 & 0.688 & 0.091 & 0.045 & 0.091 & \cellcolor{green!15}0.063 & \cellcolor{green!15}0.062 & 0.091 & 1.000 & 0.091 & 0.302 & 0.545 & 0.091 & 0.218 & 0.091 & 0.091 & 0.091 \\
  & 30 & 0.545 & 0.444 & 0.615 & 0.182 & 0.000 & 0.264 & \cellcolor{green!15}0.209 & \cellcolor{green!15}0.205 & 0.364 & 0.727 & 0.264 & 0.514 & 0.545 & 0.091 & 0.098 & 0.182 & 0.364 & 0.182 \\
  & 40 & 0.364 & 0.364 & 0.364 & 0.136 & 0.000 & 0.132 & \cellcolor{green!15}0.087 & \cellcolor{green!15}0.085 & 0.364 & 0.364 & 0.132 & 0.364 & 0.364 & -0.273 & -0.273 & 0.136 & 0.273 & 0.000 \\
  & 50 & 0.632 & 0.682 & 0.562 & 0.312 & 0.000 & 0.337 & \cellcolor{green!15}0.249 & \cellcolor{green!15}0.243 & 0.682 & 0.562 & 0.384 & 0.619 & 0.622 & 0.244 & 0.244 & 0.282 & 0.473 & 0.191 \\
  & 60 & 0.409 & 0.381 & 0.435 & 0.000 & 0.000 & 0.165 & \cellcolor{green!15}0.112 & \cellcolor{green!15}0.110 & 0.364 & 0.455 & 0.165 & 0.407 & 0.409 & -0.182 & -0.183 & 0.000 & 0.364 & 0.000 \\
\midrule
\multirow{8}{*}{\makecell[l]{\textbf{Nemo-H-47B}\\\textbf{(r-on)}}}
  & 8  & 0.668 & 0.594 & 0.720 & 0.338 & 0.002 & 0.346 & \cellcolor{green!15}0.293 & \cellcolor{green!15}0.288 & 0.485 & 0.852 & 0.413 & 0.643 & 0.668 & 0.337 & 0.362 & 0.338 & 0.356 & 0.337 \\
  & 10 & 0.641 & 0.638 & 0.645 & 0.301 & 0.001 & 0.302 & \cellcolor{green!15}0.251 & \cellcolor{green!15}0.246 & 0.631 & 0.651 & 0.411 & 0.641 & 0.641 & 0.283 & 0.283 & 0.301 & 0.321 & 0.286 \\
  & 15 & 0.582 & 0.482 & 0.650 & 0.193 & 0.000 & 0.202 & \cellcolor{green!15}0.160 & \cellcolor{green!15}0.157 & 0.388 & 0.777 & 0.301 & 0.549 & 0.583 & 0.165 & 0.179 & 0.193 & 0.227 & 0.171 \\
  & 20 & 0.556 & 0.406 & 0.646 & 0.144 & 0.000 & 0.153 & \cellcolor{green!15}0.121 & \cellcolor{green!15}0.118 & 0.303 & 0.810 & 0.245 & 0.495 & 0.556 & 0.113 & 0.131 & 0.144 & 0.173 & 0.130 \\
  & 30 & 0.537 & 0.346 & 0.641 & 0.116 & 0.000 & 0.126 & \cellcolor{green!15}0.097 & \cellcolor{green!15}0.095 & 0.245 & 0.828 & 0.203 & 0.450 & 0.537 & 0.073 & 0.090 & 0.116 & 0.152 & 0.093 \\
  & 40 & 0.505 & 0.067 & 0.663 & 0.021 & 0.000 & 0.022 & \cellcolor{green!15}0.017 & \cellcolor{green!15}0.016 & 0.035 & 0.975 & 0.034 & 0.186 & 0.505 & 0.010 & 0.029 & 0.021 & 0.028 & 0.017 \\
  & 50 & 0.503 & 0.030 & 0.666 & 0.009 & 0.000 & 0.009 & \cellcolor{green!15}0.007 & \cellcolor{green!15}0.007 & 0.015 & 0.991 & 0.015 & 0.123 & 0.503 & 0.006 & 0.027 & 0.009 & 0.010 & 0.008 \\
  & 60 & 0.500 & --    & 0.667 & 0.000 & 0.000 & 0.000 & \cellcolor{green!15}0.000 & \cellcolor{green!15}0.000 & 0.000 & 1.000 & 0.000 & 0.000 & 0.500 & 0.000 & --    & 0.000 & 0.000 & 0.000 \\
\midrule
\multirow{7}{*}{\makecell[l]{\textbf{Qwen3-14B}\\\textbf{(think)}}}
  & 8  & 0.853 & 0.829 & 0.871 & 0.706 & 0.473 & 0.708 & \cellcolor{green!15}0.638 & \cellcolor{green!15}0.629 & 0.711 & 0.995 & 0.708 & 0.841 & 0.853 & 0.706 & 0.737 & 0.706 & 0.711 & 0.706 \\
  & 10 & 0.755 & 0.680 & 0.802 & 0.517 & 0.351 & 0.514 & \cellcolor{green!15}0.445 & \cellcolor{green!15}0.436 & 0.516 & 0.998 & 0.515 & 0.717 & 0.757 & 0.513 & 0.585 & 0.515 & 0.518 & 0.515 \\
  & 15 & 0.568 & 0.478 & 0.632 & 0.311 & 0.096 & 0.297 & \cellcolor{green!15}0.233 & \cellcolor{green!15}0.227 & 0.391 & 0.749 & 0.293 & 0.541 & 0.570 & 0.140 & 0.150 & 0.311 & 0.385 & 0.177 \\
  & 20 & 0.516 & 0.604 & 0.379 & 0.225 & 0.000 & 0.212 & \cellcolor{green!15}0.157 & \cellcolor{green!15}0.153 & 0.737 & 0.295 & 0.218 & 0.467 & 0.516 & 0.033 & 0.037 & 0.225 & 0.290 & 0.093 \\
  & 30 & 0.501 & 0.660 & 0.066 & 0.030 & 0.000 & 0.034 & \cellcolor{green!15}0.020 & \cellcolor{green!15}0.020 & 0.967 & 0.035 & 0.034 & 0.185 & 0.501 & 0.003 & 0.007 & 0.030 & 0.035 & 0.028 \\
  & 40 & 0.500 & 0.657 & 0.079 & 0.035 & 0.000 & 0.041 & \cellcolor{green!15}0.023 & \cellcolor{green!15}0.023 & 0.957 & 0.043 & 0.041 & 0.203 & 0.500 & 0.000 & 0.000 & 0.035 & 0.043 & 0.023 \\
  & 50 & 0.539 & 0.693 & 0.072 & 0.039 & 0.000 & 0.038 & \cellcolor{green!15}0.021 & \cellcolor{green!15}0.021 & 0.967 & 0.039 & 0.038 & 0.194 & 0.503 & 0.006 & 0.015 & 0.039 & 0.039 & 0.026 \\
\midrule
\multirow{7}{*}{\makecell[l]{\textbf{Qwen3-32B}\\\textbf{(think)}}}
  & 8  & 0.824 & 0.788 & 0.850 & 0.647 & 0.482 & 0.648 & \cellcolor{green!15}0.576 & \cellcolor{green!15}0.566 & 0.656 & 0.991 & 0.650 & 0.806 & 0.823 & 0.647 & 0.687 & 0.645 & 0.655 & 0.645 \\
  & 10 & 0.692 & 0.564 & 0.762 & 0.389 & 0.308 & 0.393 & \cellcolor{green!15}0.325 & \cellcolor{green!15}0.318 & 0.399 & 0.985 & 0.393 & 0.627 & 0.692 & 0.384 & 0.474 & 0.389 & 0.399 & 0.384 \\
  & 15 & 0.580 & 0.370 & 0.685 & 0.225 & 0.096 & 0.224 & \cellcolor{green!15}0.175 & \cellcolor{green!15}0.171 & 0.245 & 0.917 & 0.225 & 0.475 & 0.581 & 0.163 & 0.220 & 0.223 & 0.245 & 0.173 \\
  & 20 & 0.548 & 0.475 & 0.604 & 0.284 & 0.046 & 0.272 & \cellcolor{green!15}0.210 & \cellcolor{green!15}0.205 & 0.404 & 0.696 & 0.281 & 0.530 & 0.550 & 0.100 & 0.104 & 0.278 & 0.389 & 0.141 \\
  & 30 & 0.497 & 0.586 & 0.361 & 0.214 & 0.000 & 0.199 & \cellcolor{green!15}0.144 & \cellcolor{green!15}0.141 & 0.711 & 0.283 & 0.202 & 0.449 & 0.497 & -0.005 & -0.006 & 0.214 & 0.283 & 0.059 \\
  & 40 & 0.505 & 0.667 & 0.041 & 0.021 & 0.000 & 0.021 & \cellcolor{green!15}0.012 & \cellcolor{green!15}0.012 & 0.989 & 0.021 & 0.021 & 0.145 & 0.505 & 0.011 & 0.043 & 0.021 & 0.021 & 0.016 \\
  & 50 & 0.583 & 0.722 & 0.167 & 0.091 & 0.000 & 0.091 & \cellcolor{green!15}0.052 & \cellcolor{green!15}0.052 & 1.000 & 0.091 & 0.091 & 0.302 & 0.545 & 0.091 & 0.227 & 0.091 & 0.091 & 0.091 \\
\midrule
\multirow{7}{*}{\makecell[l]{\textbf{QwQ-32B}\\\textbf{(think)}}}
  & 8  & 0.919 & 0.912 & 0.925 & 0.837 & 0.407 & 0.837 & \cellcolor{green!15}0.784 & \cellcolor{green!15}0.775 & 0.838 & 1.000 & 0.838 & 0.915 & 0.919 & 0.838 & 0.849 & 0.837 & 0.837 & 0.837 \\
  & 10 & 0.831 & 0.796 & 0.855 & 0.661 & 0.331 & 0.661 & \cellcolor{green!15}0.588 & \cellcolor{green!15}0.578 & 0.661 & 1.000 & 0.661 & 0.813 & 0.831 & 0.661 & 0.703 & 0.661 & 0.661 & 0.661 \\
  & 15 & 0.598 & 0.339 & 0.711 & 0.206 & 0.115 & 0.204 & \cellcolor{green!15}0.155 & \cellcolor{green!15}0.151 & 0.206 & 0.990 & 0.204 & 0.451 & 0.598 & 0.196 & 0.316 & 0.206 & 0.206 & 0.196 \\
  & 20 & 0.566 & 0.234 & 0.697 & 0.132 & 0.041 & 0.132 & \cellcolor{green!15}0.089 & \cellcolor{green!15}0.087 & 0.132 & 1.000 & 0.132 & 0.364 & 0.566 & 0.132 & 0.266 & 0.132 & 0.132 & 0.132 \\
  & 30 & 0.512 & 0.224 & 0.645 & 0.124 & 0.008 & 0.113 & \cellcolor{green!15}0.080 & \cellcolor{green!15}0.078 & 0.140 & 0.884 & 0.124 & 0.352 & 0.512 & 0.025 & 0.037 & 0.124 & 0.140 & 0.050 \\
  & 40 & 0.559 & 0.554 & 0.565 & 0.309 & 0.000 & 0.298 & \cellcolor{green!15}0.233 & \cellcolor{green!15}0.228 & 0.530 & 0.591 & 0.313 & 0.560 & 0.560 & 0.121 & 0.121 & 0.309 & 0.436 & 0.155 \\
  & 50 & 0.500 & 0.645 & 0.154 & 0.091 & 0.000 & 0.083 & \cellcolor{green!15}0.046 & \cellcolor{green!15}0.046 & 0.909 & 0.091 & 0.083 & 0.287 & 0.500 & 0.000 & 0.000 & 0.091 & 0.091 & 0.000 \\
\bottomrule
\end{tabular}%
}
\end{table*}

The graph 3-coloring extension reinforces the same measurement point in a second matched decision problem, now with the OpenAI run and the two NVIDIA reasoning-on runs added to the same CSA table.  The rerun applies the same JSON-tail fallback and the same pair-scenario accounting to every discovered graph-coloring model and every available $N$ row, so the comparison is now uniform across providers.  The small-$N$ frontier is strong for GPT-OSS-20B and Llama-Nemo-49B: at $N=10$, both reach ADR $=0.818$, while GPT-OSS-20B also keeps ADR$^{+w}=0.818$ and Llama-Nemo-49B keeps ADR$^{+w}=0.227$, indicating that label correctness and constructive coloring validity can diverge sharply even when separability remains high.  Nemo-H-47B behaves differently: its ordinary ADR is moderate at small $N$ ($0.338$ at $N=8$, $0.301$ at $N=10$), but ADR$^{+w}$ is almost zero from the start, showing that most classification wins do not survive assignment verification.  The three Qwen-family rows preserve the earlier pattern: ordinary separability remains nontrivial longer than constructive success, with ADR$^{+w}$ already at or near zero by $N=20$ for Qwen3-14B and much smaller than ADR for Qwen3-32B and QwQ-32B.  As in the SAT tables, the direct separability ability is the square of G-mean up to rounding, whereas $\hat{s}$ is the group-averaged Case-D point estimate and therefore need not equal the pooled product.

\subsection{Vulnerability Detection CSA}

\begin{table*}[ht]
\centering
\caption{Representative vulnerability detection CSA rows.  Balanced training supports two-sided separability; vulnerable-only training produces high positive recall but low negative recall.}
\label{tab:supp-vd-csa}
\renewcommand{\arraystretch}{1.12}
\resizebox{\textwidth}{!}{%
\begin{tabular}{llrrrrr}
\toprule
Training family & Example configuration & Scenario & $\theta_P$ & $\theta_N$ & $\CSACP$ & Interpretation \\
\midrule
Balanced & \texttt{add\_after\_ratio\_2\_v04} & 10p/g & 0.753 & 0.613 & 0.450 & strong two-sided separability \\
Balanced & \texttt{add\_after\_ratio\_1\_v01} & 1p/g & 0.754 & 0.645 & 0.447 & strong two-sided separability \\
Imbalanced & \texttt{add\_after\_percent\_50} & 10p/g & 0.913 & 0.278 & 0.245 & borderline under $\tau=0.25$ \\
Vulnerable-only & \texttt{func\_before} & 10p/g & 0.946 & 0.079 & 0.069 & one-sided vulnerable bias \\
Vulnerable-only & \texttt{func\_before\_1000\_2000} & 1p/g & 0.926 & 0.058 & 0.040 & one-sided vulnerable bias \\
\bottomrule
\end{tabular}%
}
\end{table*}

\section{Additional Interpretation}
\label{sec:supp-interpretation}

\subsection{Why CSA Is Not an Accuracy Replacement}

CSA is not meant to replace accuracy, F1, or task-specific metrics.  It answers a different question.  Accuracy asks how often the final labels match the benchmark.  CSA asks how much two-sided matched separability can be lower-bounded under a confidence protocol.  Both are useful.  A high-accuracy, low-CSA row may still be useful for engineering deployment if the deployment distribution resembles the benchmark, but it supports a weaker scientific claim about matched reasoning behavior.

\subsection{Why CSA Should Be Reported With Positive and Negative Recalls}

The product lower bound is interpretable only when readers can see both sides of the contrast.  A low CSA score can arise from weak positive-side recognition, weak negative-side recognition, or both.  Therefore, every CSA report should include $\hat\theta_P$, $\hat\theta_N$, $\hat{s}$, $\CSACP$, and the group construction rule.

\subsection{Relationship to Witness-Aware Evaluation}

Witness-aware evaluation is stronger when available.  In SAT, a satisfying assignment provides direct evidence on the SAT side.  In program repair, a test-passing patch may provide behavioral evidence.  In vulnerability analysis, an exploit, patch, or causal trace may support stronger claims.  CSA should be read as a category-only lower-bound method for settings where such richer artifacts are unavailable, incomplete, or difficult to validate at scale.

\section{Additional Threats to Validity}
\label{sec:supp-threats}

\paragraph{Mechanism identification.}
CSA certifies observable two-sided separability, not the internal mechanism used by the model.  A high CSA score is compatible with several possible mechanisms.  This is a limitation but also a design choice: the paper avoids claiming more than category outcomes can identify.

\paragraph{Matched group quality.}
The method depends on the quality of the matched groups.  If the positive and negative arms differ in unintended artifacts, CSA can still be high.  Benchmark authors should document generation rules, editing procedures, randomization, and stability checks.

\paragraph{Statistical dependence.}
The CP and Wilson calculations assume a binomial sampling model.  Generated instances may share seeds, templates, or prompts.  The group-level cluster bootstrap of \Cref{sec:supp-bootstrap} resamples matched pairs and provides the dependence-robust lower bound $\CSAboot$; we recommend reporting it beside $\CSACP$ and using it for the boundary decision whenever the two disagree.

\paragraph{Threshold choice.}
The default $\tau=0.25$ threshold is an interpretable reporting threshold, not a universal standard.  High-stakes applications should use stricter thresholds and should prefer witness-aware or causal evidence when available.

\paragraph{Domain coverage.}
SAT and vulnerability detection cover two reasoning-style settings, but additional SE tasks such as defect prediction, bug triage, flaky-test prediction, configuration reasoning, and program-repair validation would further strengthen external validity.

\section{Conclusion}

The revised supplementary material supports the main paper's new metric choice: the primary quantity is the lower-confidence bound on matched separability ability.  This removes the earlier margin-over-baseline machinery and makes the central claim cleaner.  The method no longer says that category evidence exceeds a specified non-semantic baseline.  It says that, under a matched-group construction and confidence protocol, the model has at least a stated amount of observable two-sided separability.  This is a narrower, more defensible, and more auditable claim for category-outcome reasoning evaluation.

\bibliographystyle{IEEEtran}
\bibliography{paper}

\end{document}
